\documentclass[a4paper,12pt]{article}

\usepackage[utf8]{inputenc}
\usepackage[brazil]{babel}

\title{Plano de Projeto - Brechó Online}
\author{}
\date{}

\begin{document}

\maketitle

\section{Introdução}
A compra e venda de roupas usadas muitas vezes ocorre de forma informal, por meio de redes sociais e aplicativos de mensagens, o que dificulta a organização, a busca por produtos e a segurança nas negociações. Nesse contexto, surge a necessidade de uma plataforma centralizada que facilite a interação entre compradores e vendedores, proporcionando maior eficiência, visibilidade e confiabilidade no processo.

\section{Problema}
Falta de uma plataforma centralizada que permita anunciar, buscar e negociar roupas usadas de forma simples e confiável.

\section{Objetivos}
O objetivo geral é desenvolver um marketplace para compra e venda de roupas usadas. Os objetivos específicos do projeto envolvem: 

\begin{itemize}
    \item Cadastro e autenticação de usuários
    \item Busca e filtragem de produtos
    \item Criação e gerenciamento de anúncios
    \item Comunicação entre usuários
    \item Compra e venda de produtos
\end{itemize}

\section{Possíveis Soluções}
\begin{itemize}
    \item Desenvolvimento de um site (plataforma web) acessível pelo navegador, onde usuários podem anunciar e buscar roupas usadas de forma simples e organizada.
    \item Desenvolvimento de um aplicativo mobile, focado em facilitar o acesso e a usabilidade no celular.
    \item Utilização de uma plataforma híbrida, começando com um site e evoluindo futuramente para um aplicativo, conforme a necessidade.
\end{itemize}

\section{Componentes Necessários}

\begin{itemize}
    \item \textbf{Plataforma do sistema:} site ou aplicação onde os usuários possam se cadastrar, anunciar produtos e realizar buscas.

    \item \textbf{Armazenamento:} banco de dados responsável por guardar informações como usuários, anúncios, pedidos, etc. Serviços de armazenamento de imagens e mídias.

    \item \textbf{Sistema de comunicação:} mecanismo que permita a interação entre compradores e vendedores.

    \item \textbf{Sistemas de pagamentos:} mecanismo que permita que a compra e venda dos produtos possa ser feita diretamente na plataforma.

    \item \textbf{Infraestrutura de implantação:} ambiente onde o sistema ficará disponível para acesso contínuo.

\end{itemize}

\section{Equipe}
Equipe de 6 pessoas organizada por funcionalidades:
\begin{itemize}
    \item Gestão de usuários
    \item Gestão de anúncios
    \item Busca e navegação
    \item Comunicação entre usuários
    \item Infraestrutura de pagamentos
    \item Integração, testes e análise
\end{itemize}

\section{Cronograma}

O projeto será desenvolvido ao longo de aproximadamente cinco meses, acompanhando a duração da disciplina. As atividades serão divididas em etapas principais:

\begin{itemize}
    \item \textbf{Planejamento e definição do projeto:} levantamento de requisitos, definição das funcionalidades principais e organização da equipe.

    \item \textbf{Modelagem do sistema:} definição da estrutura geral do sistema, modelagem dos dados e planejamento da arquitetura da aplicação.

    \item \textbf{Desenvolvimento das funcionalidades principais:} implementação das funcionalidades essenciais, como cadastro de usuários, criação de anúncios e sistema de busca.

    \item \textbf{Implementação de funcionalidades complementares:} desenvolvimento de comunicação entre usuários, melhorias de navegação e integração com sistema de pagamentos.

    \item \textbf{Integração e testes:} verificação do funcionamento conjunto dos componentes do sistema, identificação e correção de erros.

    \item \textbf{Finalização e entrega:} ajustes finais, documentação do projeto e preparação da apresentação.
\end{itemize}

\section{Riscos}

Durante o desenvolvimento do projeto alguns riscos podem surgir e impactar o andamento das atividades.

\begin{itemize}
    \item \textbf{Atrasos no desenvolvimento:} dificuldades técnicas ou falta de tempo podem comprometer o cronograma do projeto.

    \item \textbf{Problemas de integração:} diferentes partes do sistema podem apresentar dificuldades ao serem integradas.

    \item \textbf{Mudanças no escopo:} novas ideias ou funcionalidades podem surgir ao longo do desenvolvimento, aumentando a complexidade do sistema.

    \item \textbf{Dependência de serviços externos:} o uso de sistemas de pagamento ou serviços de armazenamento pode gerar limitações ou dificuldades de integração.
\end{itemize}

Para reduzir esses riscos, serão realizadas reuniões periódicas da equipe, acompanhamento do progresso do projeto e priorização das funcionalidades essenciais.

\section{Interfaces e Dependências}

O sistema depende da interação entre diferentes componentes internos e também de alguns serviços externos.

Internamente, os módulos responsáveis por usuários, anúncios, busca, comunicação e pagamentos precisam funcionar de forma integrada para garantir a experiência completa da plataforma.

Externamente, o sistema pode depender de serviços como plataformas de hospedagem, sistemas de pagamento online e serviços de armazenamento de imagens. Essas dependências permitem ampliar as funcionalidades da plataforma sem a necessidade de desenvolver todos os recursos internamente.

\section{Estimativas de Custos}

Por se tratar de um projeto acadêmico, os custos financeiros são relativamente baixos.

\begin{itemize}
    \item \textbf{Infraestrutura de hospedagem:} utilização de serviços de hospedagem ou servidores virtuais para disponibilizar a aplicação online.

    \item \textbf{Registro de domínio (opcional):} aquisição de um endereço próprio para acesso à plataforma.

    \item \textbf{Serviços externos:} possíveis custos associados a serviços de armazenamento ou sistemas de pagamento.

    \item \textbf{Desenvolvimento:} principal investimento relacionado ao tempo de dedicação da equipe no desenvolvimento do projeto.
\end{itemize}

Grande parte das ferramentas utilizadas no desenvolvimento possui versões gratuitas ou educacionais, o que reduz significativamente os custos do projeto.

\section{Questões em Aberto}

Algumas decisões ainda podem ser definidas ao longo do desenvolvimento do projeto:

\begin{itemize}
    \item Definição do modelo de comunicação entre compradores e vendedores.
    \item Escolha de integrar um sistema de pagamentos a plataforma, assim como a escolha de qual sistema usar.
    \item Estratégia de armazenamento das imagens dos anúncios.
    \item Possível evolução da plataforma para aplicação mobile no futuro.
\end{itemize}

Essas decisões serão tomadas pela equipe conforme o avanço do desenvolvimento e as necessidades identificadas durante o projeto.

\end{document}